% Sections 10.1 to 10.5, from lectures/L10/appendix/a_theory.md.

\section{Exponentialfunktioner och ekvationer}
\label{c10:sec:exponential}
En ekvation med en \textbf{exponentialfunktion} $f(x) = a^x$ ($a > 0$, $a \neq 1$) kan lösas för
$x$ med hjälp av logaritmer.

\textbf{Att lösa $a^x = b$:} logaritmera båda led.
\[
  x \cdot \log a = \log b \quad \Rightarrow \quad x = \frac{\log b}{\log a}
\]

\section{Logaritmer}
\label{c10:sec:logaritmer}
\textbf{Tiologaritmen} $\log_{10}$, som vanligen skrivs $\log$, definieras av
\[
  \log_{10}(10^x) = x \quad \Leftrightarrow \quad 10^{\log x} = x.
\]
Den \textbf{naturliga logaritmen} $\ln$ har basen $e \approx 2{,}718$:
\[
  \ln(e^x) = x \quad \Leftrightarrow \quad e^{\ln x} = x
\]

\section{Logaritmlagar}
\label{c10:sec:logaritmlagar}
För logaritmer gäller följande räknelagar:
\begin{gather*}
  \log(ab) = \log a + \log b \\
  \log\left(\frac{a}{b}\right) = \log a - \log b \\
  \log(a^n) = n \cdot \log a
\end{gather*}
\textbf{Basbyte:}
\[
  \log_a b = \frac{\log b}{\log a} = \frac{\ln b}{\ln a}
\]

\section{Decibel (dB)}
\label{c10:sec:decibel}
Decibel används inom elektroteknik för att ange förstärkning och dämpning på en logaritmisk skala.

\textbf{Spänningsförstärkning:}
\[
  G_{\text{dB}} = 20 \log_{10}\left(\frac{U_{\text{ut}}}{U_{\text{in}}}\right)
\]
Omvänt:
\[
  \frac{U_{\text{ut}}}{U_{\text{in}}} = 10^{G_{\text{dB}}/20}
\]
\textbf{Effektförstärkning:}
\[
  G_{\text{dB}} = 10 \log_{10}\left(\frac{P_{\text{ut}}}{P_{\text{in}}}\right)
\]
\textbf{Nivå i dBV}, decibel relativt $1\,\text{V}$:
\[
  U_{\text{dBV}} = 20 \log_{10}\left(\frac{U_{\text{RMS}}}{1\,\text{V}}\right)
\]

\needspace{12\baselineskip}
\section{Typexempel}
\label{c10:sec:typexempel}

\begin{exempel}[Lösa logaritmiska ekvationer]
Lös \textbf{a)} $3^x = 81$, \textbf{b)} $2^{x-1} = 64$ och \textbf{c)} $e^{x-2} = 150$.

\textbf{a)} $3^x = 3^4 \Rightarrow x = 4$. Alternativt är $x = \frac{\log 81}{\log 3} = 4$.

\textbf{b)}
\[
  (x-1)\log 2 = \log 64 \quad \Rightarrow \quad x - 1 = \frac{\log 64}{\log 2} = 6
  \quad \Rightarrow \quad x = 7
\]

\textbf{c)}
\[
  x - 2 = \ln 150 \approx 5{,}01 \quad \Rightarrow \quad x \approx 7{,}01
\]
\end{exempel}

\begin{exempel}[Halveringstid för batteriladdning]
Ett batteri tappar halva sin laddning på $30$ timmar:
\[
  u(30) = U_0 \cdot a^{30} = 0{,}5\,U_0
\]
\textbf{Beräkna förändringsfaktorn $a$:}
\[
  a^{30} = 0{,}5 \quad \Rightarrow \quad a = 0{,}5^{1/30} \approx 0{,}977
\]
\textbf{När återstår $20\,\%$?}
\[
  a^t = 0{,}2 \quad \Rightarrow \quad t = \frac{\log 0{,}2}{\log a}
  = \frac{\log 0{,}2}{\log 0{,}5^{1/30}} \approx 69{,}7\,\text{h}
\]
Räkna med det oavrundade värdet på $a$; med $0{,}977$ blir svaret $69{,}2\,\text{h}$.
\end{exempel}

\begin{exempel}[Linjär spänningsförstärkning]
Två signaler har nivåerna $L_1 = 20\,\text{dB}$ och $L_2 = 46\,\text{dB}$.
\begin{gather*}
  G_{\text{dB}} = L_2 - L_1 = 26\,\text{dB} \\
  G_{\text{lin}} = 10^{26/20} = 10^{1{,}3} \approx 20
\end{gather*}
\end{exempel}

\begin{exempel}[dBV till volt]
En sinusspänning har nivån $31{,}0\,\text{dBV}$.
\begin{gather*}
  U_{\text{RMS}} = 10^{31{,}0/20} \approx 35{,}5\,\text{V} \\
  \abs{U} = U_{\text{RMS}} \cdot \sqrt{2} \approx 35{,}5 \cdot 1{,}414 \approx 50{,}2\,\text{V}
\end{gather*}
\end{exempel}
